
\documentclass[template=tabling,81pt,headonall]{azmoon}
\usepackage{xepersian}
\usepackage{amsfonts}
\usepackage{graphicx}
\usepackage{svg}
\svgpath{ {./images/} }
\graphicspath{ {./images/} }
\settextfont{Yas}
\setdigitfont{A Iranian Sans}
\usepackage{fontawesome5}

\printanswers
    \teacher{محمد صالح علی اکبری}
    \teachertitle{دبیر}
    \city{گناباد}
    \schooltitle{دبیرستان}
    \school{شهید جوانبخت}
    \grade{هفتم (نمونه ۲)}
    \branch{نمونه ۲}
    \topic{ریاضی}
    \examdate{۲۵ مهر ۱۴۰۵}
    \answertime{۳۵ دقیقه}
    \begin{document}
	\begin{questions}
		\nointerlineskip%
		\vskip-\baselineskip
		\question{%
یک ساختمان $5$ طبقه ساخته شده است. بالاترین طبقهٔ آن $2$ واحد دارد، اما هر یک طبقه که پایین می‌رویم، $2$ واحد به تعداد واحدها اضافه می‌شود.\\
الف) چه الگویی بین تعداد واحدها در طبقه‌های این ساختمان وجود دارد؟\\
ب) این ساختمان در کل چند واحد دارد؟‌
\\‌
\\‌
\\‌
\\}\question{%
با انگشتان یک دست به $5$ صورت می‌توان عدد $1$ را نشان داد. به چند صورت می‌توان عدد $2$ را نشان داد؟ (تمام حالت‌های مختلف و روش خود را بنویسید)‌
\\‌
\\‌
\\‌
\\}\question{%
$20$ دستگاه دوچرخه و سه‌چرخه در یک پارکینگ وجود دارد. اگر تعداد کل چرخ‌های آن‌ها $45$ عدد باشد، چند دوچرخه و چند سه‌چرخه در توقفگاه وجود دارد؟ (راه‌حل یا جدول حدس و آزمایش را کامل بنویسید)‌
\\‌
\\‌
\\‌
\\‌
\\‌
\\‌
\\‌
\\‌
\\}\end{questions}
    \end{document}
    